\titleformat{\chapter}[display]
  {\normalfont\bfseries}{}{0pt}{\Large} % Este comando es para eliminar la aparición de "Capítulo N" en el PDF. Para emplear el formato de la tesis simplemente eliminamos la línea. 
\chapter{Introducción}
El virus del dengue (DENV) es un miembro de la familia Flaviviridae transmitido principalmente por la picadura del mosquito hembra Aedes aegypti. Este virus se clasifica en cuatro serotipos distintos (DENV-1, DENV-2, DENV-3 y DENV-4) de acuerdo con su complejo antigénico de reconocimiento, los cuales comparten entre un 65\% y un 70\% de similitud en su secuencia de aminoácidos. A nivel mundial, el DENV causa entre 50 y 100 millones de infecciones reportadas al año; una incidencia fuertemente influenciada por factores de riesgo como la urbanización acelerada, la presencia de canales o agua estancada, y el cambio climático, los cuales contribuyen directamente a la replicación y dispersión del vector. En el caso específico de México, los reportes epidemiológicos correspondientes al año 2026 continúan demostrando una alta incidencia anual y una tasa significativa de defunciones asociadas en el país.
%(1) Aspectos genéticos del virus del dengue, Genetic aspects of dengue virus, Mike Contreras1* María Angélica Rincón2, Andrea Vásquez Guillén3 Ruth Moreira4, Diana Callejas5. 
El genoma del DENV consiste en una sola cadena de ARN monocatenario de sentido positivo (+ssRNA) de aproximadamente 10.7 a 11 kb. Este genoma codifica un único marco abierto de lectura (ORF) compuesto por cerca de 3400 codones, el cual se traduce en una poliproteína que posteriormente es procesada en tres proteínas estructurales: cápside (C), de membrana (prM) y envoltura (E) y siete proteínas no estructurales: NS1, NS2A, NS2B, NS3, NS4A/2K, NS4B y NS5. Específicamente, el complejo proteico NS2B-NS3 (proteasa) juega un rol clave tanto en el procesamiento estructural de los componentes del virión como en la replicación viral. No obstante, es importante destacar que, de todos los serotipos de este virus, el serotipo 2 (DENV-2) es el más estudiado debido a su alta prevalencia y por ser el mayor causante de manifestaciones graves como la fiebre por dengue y el dengue hemorrágico \cite{contreras2021aspectos}. Desafortunadamente, a la fecha no se cuenta con vacunas con una eficacia global óptima o tratamientos antivirales aprobados universalmente, por lo que el desarrollo de inhibidores dirigidos contra las proteínas responsables de la replicación de DENV-2 es una prioridad. Entre los compuestos evaluados, se han reportado inhibidores ortostéricos (como derivados de quinoxalina) e inhibidores alostéricos (como derivados de benzofurano) con una buena concentración inhibitoria in vitro y un gran impacto sobre sitios específicos; sin embargo, limitaciones farmacocinéticas o de toxicidad han impedido que avancen hacia fases clínicas. Por lo tanto, el diseño de nuevos inhibidores contra la proteasa NS2B-NS3 (utilizando como diana la estructura de la proteína en conformación inactiva-abierta, identificada con el código PDB: 2FOM) representa una estrategia terapéutica prometedora debido a que al estabilizar el estado inactivo-abierto se impide que la enzima realice el cambio conformacional hacia su forma cerrada y catalíticamente activa, bloqueando así el procesamiento de la poliproteína viral. 

Como alternativa en el desarrollo de fármacos, en este estudio se reporta el análisis de nanoestructuras de carbono. Debido a la versatilidad de los átomos de carbono para adoptar diferentes hibridaciones orbitales, estos pueden formar enlaces direccionados sumamente estables y dar origen a diversas estructuras alotrópicas. Una de las más destacadas es el fulereno, una estructura molecular compuesta únicamente por carbono que se presenta en distintos tamaños, como C\textsubscript{20}, C\textsubscript{60}, C\textsubscript{70} y C\textsubscript{76}. La más conocida y estudiada de estas formas es el C\textsubscript{60} (o buckminsterfulereno), constituido por 20 hexágonos y 12 pentágonos. Esta molécula cuenta con dos tipos de enlaces: los enlaces [6,6] con una distancia de 1.38 Å, y los enlaces [5,6] con una distancia de 1.45 Å. \cite{fajardo2012dengue} El descubrimiento de estas nuevas formas de carbono otorgó a los investigadores Harold W. Kroto, Robert F. Curl Jr. y Richard E. Smalley el Premio Nobel de Química en 1996.

A pesar de sus notables propiedades físicas y electrónicas, el C\textsubscript{60} presenta una limitación crítica para su uso biomédico: su baja solubilidad en medios acuosos. Para superar este inconveniente, se ha propuesto la funcionalización del fulereno mediante la adición de aminoácidos. Numerosos estudios de aductos de C\textsubscript{60} modificados con L-metionina, L-hidroxiprolina y L-arginina han demostrado una pronunciada actividad biológica como agentes antioxidantes, antimicrobianos, antivirales, antitumorales, neuroprotectores, membranotrópicos y fotodinámicos, actuando incluso como radioprotectores e inductores de la apoptosis. Específicamente, el aducto de C\textsubscript{60} funcionalizado con cuatro moléculas de glicina ha sido reportado con una notable actividad antioxidante, propiedades de fotoprotección y una excelente hemocompatibilidad, logrando además el parámetro más importante: una solubilidad aceptable para aplicaciones en sistemas biológicos \cite{krishna2024dengue}. Así, estos aductos sirven como andamios moleculares versátiles para el desarrollo de nuevos materiales con fines médicos y terapéuticos contra diversas enfermedades.

Dado lo anterior, el presente estudio teórico-computacional consiste en la determinación de una serie de parámetros estructurales y energéticos que forman parte de la caracterización del aducto C\textsubscript{60}-glicina (modificado con cuatro moléculas de glicina). Mediante una metodología computacional, este trabajo tiene como objetivo evaluar y analizar el comportamiento, la afinidad y la actividad de este aducto como un inhibidor potencial contra la proteasa NS2B-NS3 del virus del dengue serotipo 2 (DENV-2), empleando para ello el modelo estructural en conformación abierta-inactiva (PDB: 2FOM), debido a que esta conformación impide que la proteína adopte su estructura activa-cerrada necesaria para llevar a cabo la catálisis proteolítica. 



